\begin{Q}
\textbf{CLIP-Score Guided LoRA Fine-Tuning for Stable Diffusion 1.5. 50 Points}

In this problem, you will adapt pretrained \emph{Stable Diffusion 1.5} with a small LoRA module.
The target concept is a single prompt family centered on
\[
\texttt{a Monet painting of a flying pig}.
\]
The training setup is intentionally simple:
\begin{itemize}
\item create a bank of 16 generated training prompts,
\item for each prompt, generate 8 candidate images,
\item score the 8 candidates with a frozen CLIP-style reward model,
\item keep the top-1 image,
\item do one LoRA-only diffusion fine-tuning update on that selected image.
\end{itemize}
So the visible training loop is ``sample 8, pick top 1, update'', repeated over the 16 prompt variants.

The visible files are
\[
\texttt{data/train.jsonl},\quad
\texttt{data/val.jsonl}.
\]
The hidden file \texttt{data/test.jsonl} will \emph{not} be given to students; it will be used by the autograder after the LoRA checkpoint is saved.
The visible training split contains exactly 16 generated training prompts.
The visible validation split contains exactly 1 prompt.
These train/val/test prompts come from the data-generation file, and the hidden test prompt is not part of the 16-prompt training bank.
Students should run the script only in \texttt{--mode train} locally.

You are provided with a starter Python file.
Your task is to fill in \emph{only} the functions whose docstring begins with \textbf{Implement:}.
For clarity, the main graded functions in this assignment are:
\begin{align*}
&\texttt{build\_prompt\_bank},\quad
\texttt{generate\_candidate\_images},\\
&\texttt{pick\_top1\_candidate},\quad
\texttt{train\_lora\_sft}.
\end{align*}

\medskip
\noindent\textbf{Training protocol.}
The released code follows this procedure:
\begin{itemize}
\item load Stable Diffusion 1.5 from \texttt{diffusers},
\item freeze the base model and attach a trainable UNet LoRA adapter,
\item iterate over the 16 visible training prompts,
\item for each prompt, generate 8 candidate images with at least 20 diffusion inference steps and compute CLIP scores,
\item keep only the highest-scoring image,
\item update only the LoRA weights with a diffusion training step on the selected image,
\item evaluate on one visible validation prompt by generating 4 images,
\item save a checkpoint containing \emph{LoRA weights only}.
\end{itemize}

\medskip
\noindent\textbf{Validation report.}
For the visible validation prompt, students must show:
\begin{itemize}
\item 4 generated images,
\item the CLIP score for each image,
\item the average score without LoRA,
\item the average score with LoRA,
\item the improvement from LoRA.
\end{itemize}
The code saves a validation image sheet for this purpose.

\medskip
\noindent\textbf{Hidden evaluation.}
Gradescope will load the submitted LoRA-only checkpoint and evaluate it on one hidden test prompt.
For transparency, the released script reports both the frozen-base average score and the LoRA-adapted average score on 4 generated test images, together with the improvement.
Students do not see the hidden test prompt.

\medskip
\noindent\textbf{Autograder.}
Gradescope will:
\begin{itemize}
\item run unit tests on the four core functions above,
\item run your training script,
\item load \texttt{outputs/diffusion\_model.pt},
\item verify that the checkpoint contains LoRA weights only,
\item evaluate the saved adapter on the hidden test prompt.
\end{itemize}
The score breakdown is:
\begin{itemize}
\item 40 points for correct implementations of the four core functions,
\item 10 points for reporting the required validation generations and average score from \texttt{--mode train}.
\end{itemize}

\begin{enumerate}

\item \textbf{\texttt{build\_prompt\_bank}: construct the 16 training prompts (10 points).}\\
Implement \texttt{build\_prompt\_bank()}.
It should return exactly the 16 training prompts generated by the data-generation script.
These are held-in training prompts related to the overall image-generation task, but they are not copies of the held-out validation or hidden-test prompt.

\item \textbf{\texttt{generate\_candidate\_images}: generate 8 images for one prompt (10 points).}\\
Implement \texttt{generate\_candidate\_images(...)}.
For one prompt and a list of seeds, it should deterministically produce the candidate images with Stable Diffusion 1.5.
Use at least 20 inference steps for every generated image.

\item \textbf{\texttt{pick\_top1\_candidate}: select the best image by CLIP score (10 points).}\\
Implement \texttt{pick\_top1\_candidate(...)}.
It should score all generated candidates with frozen CLIP, identify the top-1 image, and return the selected image, its score, and its seed.

\item \textbf{\texttt{train\_lora\_sft}: LoRA fine-tuning with top-1 self-training (10 points).}\\
Implement \texttt{train\_lora\_sft(...)}.
The training loop should:
\begin{itemize}
\item iterate over the 16 visible training prompts,
\item generate 8 candidates per prompt with at least 20 inference steps,
\item pick the top-1 candidate by CLIP score,
\item do one LoRA-only diffusion update using that selected image,
\item evaluate after each epoch on the visible validation prompt,
\item keep the best LoRA weights by validation average score and restore them before saving.
\end{itemize}

\item \textbf{Validation-Mode Report and Checkpoint Submission (10 points).}\\
Run the script in \texttt{--mode train} and report:
\begin{itemize}
\item the visible validation prompt,
\item the 4 per-image validation scores without LoRA,
\item the 4 per-image validation scores with LoRA,
\item the average validation score without LoRA,
\item the average validation score with LoRA,
\item the improvement from LoRA,
\item the saved validation image sheet.
\end{itemize}
Students do not have access to the hidden test prompt.
Submit \texttt{outputs/diffusion\_model.pt} to Gradescope, and the autograder will evaluate it on hidden test data, reporting both the frozen-base and LoRA-adapted average scores over 4 images.

\begin{solution}

Reference numbers depend on the installed SD 1.5 checkpoint, CLIP checkpoint, GPU, and runtime environment.
The key requirement is that the implementation uses Stable Diffusion 1.5, frozen CLIP scoring, UNet LoRA, and at least 20 inference steps per generated image.

\bigskip

\begin{tabular}{lccc}
\hline
Split & Base avg & LoRA avg & Improvement \\
\hline
Validation & -- & -- & -- \\
Hidden test & -- & -- & -- \\
\hline
\end{tabular}

\bigskip

Visible validation prompt: \texttt{an impressionist Monet style painting of a pig flying through the sky}

\smallskip

Validation base image scores:
\[
--,\quad --,\quad --,\quad --
\]

\smallskip

Validation LoRA image scores:
\[
--,\quad --,\quad --,\quad --
\]

The checkpoint should contain LoRA weights only.

\end{solution}

\end{enumerate}
\end{Q}
